\section{Codizes}

\subsection{Adeptus Sororitas}

\begin{itemize}

 \item Eine Armee des Adeptus Sororitas benutzt den Eintrag des Adeptus
  Sororitas der Alliertenmatrix aus dem Regelbuch.

\end{itemize}

\subsection{Blood Angels}

\begin{itemize}

 \item Ein Furioso-Scriptor-Cybot zählt als ein Blood Angels Scriptor in Bezug
  auf das Generieren von Psi-Kräften aus dem RB. Beachte, dass ein
  Furioso-Scriptor-Cybot Meisterschaftsgrad 1 ist, obwohl er 2 Psi-Kräfte hat.

 \item Die Psi-Kraft Blutspeer hat keinen Effekt auf Flieger oder fliegende
  monströse Kreaturen im Flugmodus.

\end{itemize}

\subsection{Chaos Dämonen}

\begin{itemize}

 \item Wenn auf der Warpsturmtabelle ein Ergebnis ermittelt wurde, bei dem ein
  W6 für jede nicht im Nahkampf gebundene gegnerische Einheit (sowie für Dämonen
  bestimmter Gottheiten) gewürfelt werden muss, wird nicht für Einheiten in
  Transportfahrzeugen oder in Gebäuden gewürfelt.

 \item Das Ergebnis ``Dämonische Besessenheit'' von der Warpsturmtabelle hat
  keinen Effekt auf Psioniker, welche sich in Transportfahrzeugen oder in
  Gebäuden befinden. Fliegende monströse Kreaturen im Jagdflug (die gleichzeitig
  Psioniker sind) werden normal betroffen.

 \item Wenn eine neue Einheit Dämonen durch die Portalglyphe oder die
  Warpsturmtabelle generiert wird, darf der Dämonen Spieler die Anzahl der
  Dämonen auswürfeln, bevor er sich für die Art der Dämonen entscheidet. 

 \item Wenn während des Spiels ein neuer Psioniker (oder eine Einheit Psioniker)
  generiert wird, würfle sofort die zufällig bestimmten Psi-Kräfte aus, bevor du
  die Modelle auf dem Spielfeld platzierst. Psioniker, die auf diese Weise das
  Spielfeld betreten, besitzen Warpenergie entsprechend ihrem
  Meisterschaftsgrad.

 \item Unabhängige Charaktermodelle aus dem Codex: Chaosdämonen dürfen sich
  nicht Einheiten aus dem Codex: Chaos Space Marines anschließen.

 \item Verletzungen ignorieren kann gegen Wunden, die durch dämonische
  Instabilität hervorgerufen, genutzt werden.

 \item Kairos Schicksalswebers Wiederholungswurf (den er durch den Stab des
  Morgen erhält), kann nicht genutzt werden, um den Wurf um die 6. Und 7. Runde
  zu wiederholen. Ebenso kann der Wurf zur Ermittlung des Nachtkampfes sowie
  Würfe des Gegners nicht mithilfe des Stabes neu gewürfelt werden.

 \item Kairos Schicksalsweber wird bezüglich Sonderregeln oder Angriffen, die
  Psioniker betreffen, immer als einzelner Psioniker behandelt.
 \item Kairos Schicksalswebers Wiederholungswurf, der ihm durch den Stab des
  Morgen gewährt wird und die Kriegsherrenfähigkeit Fürst der Unrealität können
  auch genutzt werden, wenn sich Kairos in Reserve befindet. 
 
 \item Wenn Kreischer ihre normalen Nahkampfattacken gegen eine einzelne
  Spezial-Attacke mit Warpmaul ersetzen, dann verfällt in diesem Falle auch die
  Bonusattacke, für den Nahkampfangriff. Niederschmettern ist hiervon nicht
  betroffen.

 \item Deckungswürfe gegen Tiefflugangriffe der Kreischer sind nur möglich, wenn
  sich die betroffenen Modelle in Geländezonen befinden (Fahrzeuge
  ausgeschlossen) oder wenn eine Sonderregel dies erlaubt.

 \item Das Blaue und Rosa Feuer des Tzeentch sind Dämonengeschenke, die zum
  erhabenen Feuerdämon gehören, nicht zum brennenden Streitwagen. Das bedeutet,
  wenn sich der Streitwagen bewegt, kann in der Schussphase das blaue Feuer nur
  als Schnellschuss und das Rosa Feuer nicht abgefeuert werden.

 \item Wenn die Maske des Slaanesh den Tanz des Fesselns gegen eine Einheit
  einsetzt, wird für jede Bewegung, die die Einheit macht, einzeln gewürfelt.
  Das bedeutet, bewegt sich die Einheit, würfle einen W3. Rennt die Einheit,
  würfle einen W3 usw. Es wird bei jeder Bewegung neu gewürfelt und der Tanz
  wirkt sich auf jede Art von Bewegung aus, auch Turbo-Boost und Vollgas.
  Fliegende monströse Kreaturen, die durch den Tanz des Fesselns betroffen sind,
  können sich nicht im Jagdflug bewegen.

 \item Der verstörende Gesang einer Slaanesh Bestie verursacht -1 auf den
  Moralwert, für jede Bestie die sich in 12 Zoll Umkreis um einen Psioniker
  befindet. Dieser Modifikator ist kumulativ mit anderen Bestien/Einheiten, die
  über die Sonderregel verstörender Gesang verfügen. Beachte, das ein Moralwert
  Test nicht das Gleiche ist wie ein Profilwerttest, daher ist es mit einer
  Doppel 1 immer möglich, einen Moralwerttest zu bestehen, selbst wenn
  Slaaneshbestien den MW auf 0 gesenkt haben.

 \item
Moralwertabzüge durch den Verdammnisstein sind dauerhaft, selbst wenn das
betroffene Modell zwischenzeitlich stirbt und durch eine Sonderregel wieder ins
Spiel kommt. Wenn der Verdammnisstein den Moralwert eines Modells auf 0 senkt,
wird es aus dem Spiel entfernt. Dagegen hilft keine Sonderregel, das Modell ist
dauerhaft aus dem Spiel entfernt.

 \item Wenn ein Modell mit mehreren Lebenspunkten einen Profilwertest, ausgelöst
  durch Ekelerregende Heimsuchung oder Pavane des Slaanesh, nicht besteht und
  dann noch lebt, endet die Psionische Attacke sofort.

 \item Jede nicht verhinderte Verwundung durch die Sonderregel Geisträuber oder
  Überträger zwingt das betroffene Modell(solange es noch lebt) zu einem
  erneuten Initiative- /Widerstandstest.

 \item Wenn eine Einheit mit Ikone des Chaos aus der Reserve erscheint, können
  nachfolgende Einheiten mit der Sonderregel Schocktruppen, diese Ikone sofort
  nutzen (das heißt im selben Spielzug), um die Abweichung beim Schocken zu
  reduzieren.

 \item Wenn ein Modell, das über die erhabene Belohnung Warpband verfügt, sich
  auf einem lahmgelegten Streitwagenbefindet und getötet wird, entferne das
  komplette Modell inklusive Streitwagen und platziere es in der aktiven
  Reserve. Das Modell muss dann das Spielfeld mittels Schocken betreten.

 \item Eine Portalglyphe kann bereits in dem Spielzug, in dem sie platziert,
  neue Einheiten generieren.

 \item Eine Portalglyphe hat keine bestimmte Höhe und kann durch die kleine
  Explosionsschablone dargestellt werden.

 \item Wenn ein Modell mit mutierender Warpklinge durch
  Gedankenkontrollskarabäen übernommen wird und in diesem Zustand ein
  Charaktermodell oder eine monströse Kreatur aus der eigenen Armee tötet, wird
  die Warp Mutation Sonderregel nicht angewendet.

 \item Wenn ein Modell, das durch Gedankenkontrollskarabäen übernommen wurde,
  ein Modell aus der eigenen Armee angreift (inklusive sich selbst) und dieses
  Modell hat die Sonderregel Brennendes Blut, kommt diese Regel nicht zur
  Anwendung.

 \item Der Bonus von +1 auf Verletzungen ignorieren Würfe durch die Sonderregel
  Warpflammen ist kumulativ mit anderen Quellen von Warpflammen. Beachte das der
  Verletzungen ignorieren Wurf niemals besser als 2+ sein kann.

 \item Wenn die dämonischen Belohnungen ausgewürfelt werden, werden die
  warpgeschmiedete Rüstung und die undurchdringliche Haut nicht als identisch
  betrachtet, bezüglich der Regel, dass ein Modell jede Belohnung nur einmal
  haben kann.

 \item Ein Modell mit der mutierenden Warpklinge muss die Waffe im Nahkampf
  genutzt haben, damit die Sonderregel der Waffe zur Anwendung kommen kann.

 \item Ein Modell mit der Blutklinge muss die Waffe im Kampf nutzen, damit die
  Sonderregel der Waffe zur Anwendung kommen kann.

 \item Die Sonderregeln Warpverderben und Warpmutation können im selben
  Initiativeschritt gegen ein und dasselbe Modell wirken.

\end{itemize}

\subsection{Chaos Space Marines}

\begin{itemize}

 \item Ein Modell ohne die Sonderregel ``Champion des Chaos'' darf heroische
  Interventionen durchführen, um ein Model mit der Regel ``Champion des Chaos''
  in einer Herausforderung abzulösen. (RB. Seite 65)

 \item Zwei unabhängige Charaktermodelle mit verschiedenen Malen des Chaos
  können sich gleichzeitig derselben Einheit ohne irgendein Mal des Chaos
  anschließen.

 \item Psi-Waffen können Kharn den Verräter nicht sofort ausschalten, auch nicht
  wenn deren Stärke doppelt so hoch wie Kharn's Widerstand ist. (Seite 59, CSM
  Codex)

 \item Ein Kriegsschmied kann kein feindliches Fahrzeug verfluchen, während er
  in einem Transportfahrzeug ist, nicht mal wenn das Fahrzeug eine Feuerluke
  hat. (Seite 34, CSM Codex)

 \item Ein Modell, das einen durch den schwarzen Streitkolben verursachten
  Widerstandstest verpatzt, wird aus dem Spiel entfernt, selbst wenn es über die
  Sonderregel ``Ewiger Krieger'' verfügt. (Seite 69, CSM Codex)

 \item Wenn ein im Nahkampf gebundenes CSM Modell am Ende der Nahkampfphase,
  innerhalb von 3 Zoll keinen Gegner mehr hat, weil diese den durch den
  schwarzen Streitkolben verursachten Widerstandstest verpatzt haben, gilt das
  Modell dennoch solange als im Nahkampf gebunden, bis alle gegnerischen
  Modelle, die an diesem Nahkampf beteiligt sind, entfernt wurden. (Seite 69,
  Codex CSM)

 \item Wenn Typhus seinen ``Vernichterschwarm'' entfesselt, während er in einer
  Herausforderung kämpft, platziere die entsprechende Schablone ganz normal über
  Typhus. Das Modell, das mit Typhus in der Herausforderung kämpft, gilt als
  eigene Einheit, unabhängig von anderen Einheiten/Modellen die von der
  Schablone betroffen sind. Verwundungen werden dann normal zugeteilt, für jede
  von der Schablone betroffene Einheit. Verluste werden normal entfernt,
  beginnend mit den Modellen, die Typhus am nächsten sind. Der
  ``Vernichterschwarm'' kann so auch Modelle verwunden, die nicht in der
  Herausforderung kämpfen. (Seite 61, Codex CSM)

 \item Seuchenzombies können Waffenstellungen besetzen und deren Waffen
  abfeuern.  (Seite 105, RB.)

 \item Die Sichtlinie für den Höllendrachen wird von seinem Kopf gezogen. Der
  Kopf hat einen 360 Grad Sichtbereich und ignoriert den Rest des Modells, wenn
  die Sichtlinie gezogen wird. (CSM FAQ)

 \item Der Dimensionschlüssel kann nicht den Effekt des Störsenders, der zu
  einem Land Speeder Storm gekauft wurde, außer Kraft setzen. (Seite 69, Codex
  CSM) 

\end{itemize}

\subsection{Dark Angels}

\begin{itemize}

 \item Eine Einheit, die Belial enthält, weicht nicht ab, wenn sie die
  Psi-Fähigkeit „Tor zur Unendlichkeit" benutzt. Gleichzeitig kann kein Modell
  der Einheit vom Warp verschlungen werden.

 \item Ezekiel erhält durch das Buch der Erlösung kein KG 6.

 \item Der Landspeeder Darkshroud verhüllt auch sich durch den Mantel der Engel
  selbst, d.h. er erhält Schleier und Tarnung.

 \item Der Widerstandmalus von -1 durch den Ravenwing-Granatwerfer wirkt sofort.
  Beispielsweise verwunden andere Schüsse der Einheit, die den Granatwerfer
  benutzt hat, besser.

\end{itemize}

\subsection{Dark Eldar}

\begin{itemize}

 \item Gegen Treffer, die durch die Sonderregel Jagdklingen der Hymeshin
  verursacht wurden, sind Deckungswürfe nur durch Geländezonen oder durch
  Sonderregeln (z.B.  Schleier, Tarnung) gestattet.

\end{itemize}

\subsection{Eldar}

\begin{itemize}

 \item Wenn ein Modell, das über die erhabene Belohnung Warpband verfügt, sich
  auf einem lahmgelegten Streitwagen befindet und getötet wird, entferne das
  komplette Modell inklusive Streitwagen und platziere es in der aktiven
  Reserve. Das Modell kann muss dann das Spielfeld mittels Schocken betreten.

 \item Die Feuerhagel-Exarchen-Fähigkeit wirkt auf jede Schusswaffe des
  Exarchen, also auch Granaten.

 \item Gegen Prinz Yriels Auge des Zorns Attacken dürfen keine Deckungswürfe
  durchgeführt werden

 \item Illic Nachtspeer muss beim Platzieren den 1“-Mindestabstand zu
  feindlichen Modellen einhalten.

 \item Wenn Illic einer andere infiltrierenden Einheit angeschlossen wird, kann
  Illics Fähigkeit Wandler auf dem verborgenem Pfad nicht eingesetzt werden. Es
  gelten die normalen Regeln zum Platzieren von Infiltratoren.

 \item Ein Spieler muss vor dem Abfeuern des Impulslasers bestimmen, welche
  Waffen das Modell abfeuern wird.

 \item Ein Verwundungswurf von 6 mit der Sonderregel Shurikensturm verwundet
  automatisch.

 \item Die Sonderregel Ewiger Krieger wirkt nicht, wenn der Moraltest verpatzt
  wurde, der durch eine Wunde des Asuryans Schwert verursacht wurde.

 \item 
  \begin{minipage}[t]{0.65\textwidth}
Das Serpentschild ist ein Ausrüstungsgegenstand. Als Mündung zählt eine der
vorderen Spitzen des Serpents. Als Waffenlauf wird die Innenseite des Rumpfes
genommen. 
\end{minipage} \
  \begin{minipage}[t]{0.35\textwidth}
\includegraphics[height=2cm, valign=T, margin=0.0cm]{Bilder/serpent_transparent.png}
\end{minipage}

 \item Wenn ein Rettungswurf mit einem Streuschild bestanden wird, so macht
  jedes Modell in 6“-Umkreis (auch der Träger des Schilds) einen Blendentest.
  Beachte, dass der Träger des Schilds den Blendentest nicht automatisch
  besteht. 

 \item Jedes Reliquien des alten Ruhms darf nur einmal enthalten sein, egal wie
  viele Haupt- oder Verbündetenkontingente eine Armee hat.

 \item Wunden durch das Mentale Duell können durch ``Achtung, Sir!'' weitergegeben
  werden. Der KG- \& BF-Malus bleibt bei dem Zielmodell.

 \item Runenleser, die vom Rat der Seher abgekoppelt werden, tun dies
  unmittelbar nachdem die Kriegsherrenfähigkeit ausgewürfelt wurde, aber bevor
  die Psikräfte generiert werden.

\end{itemize}

\subsection{Iyanden Supplement}

\begin{itemize}

 \item Ein Iyanden-Supplement verwendet die Alliierten-Matrix des Codex: Eldar.
  Zusätzlich zählen Kontingent aus dem Codex: Eldar und ein Kontingent des
  Iyanden-Supplement als Waffenbrüder.

 \item Zu Einheiten, die von der Psikraft ``Voice of Twilight'' profitieren
  wollen, muss der Psioniker eine Sichtlinie ziehen können und sich in 12“ zum
  Ziel befinden.

\end{itemize}

\subsection{Grey Knights}

\begin{itemize}

 \item Der Neuralschredder nutzt zum Verwunden den Mehrheitswiderstand.

 \item  Für die Psikraft Exzension gelten folgende Sektionen aus dem Codex:
  Chaosdämonen (S. 61–69) als Dämonengeschenke: Schusswaffen, Nahkampfwaffen,
  dämonische Eigenschaften, Ikonen \& Instrumente, Höllengeschmiedete Artefakte,
  dämonische Gaben, Dämonische Präsenzen und Dämonische Psi-Kräfte. Zusätzlich
  zählt alles, was in irgendeinem Codex/Armeeliste als Dämonengeschenk gilt.
  Beachte, dass Dämonische Reittiere und Regelbuch-Psi-Kräfte nicht von
  Exzension betroffen werden.

 \item Einheiten, die von den Psi-Kräften ``Reinigende Flamme'' und/oder
  ``Heldentod'' betroffen sind, können jeweils einen Psiresitenzwurf
  durchführen.

 \item Wenn ein Modell ``Reinigende Flamme'' wirkt, während es in einer
  Herausforderung kämpft, wird die gesamte gegnerische Einheit getroffen, nicht
  nur der Charakter, der in der Herausforderung kämpft.

\end{itemize}

\subsection{Imperiale Armee}

\begin{itemize}

 \item Einzelne Teile eines Infantriezuges können in Reserve versetzt werden,
  während der Rest des Zuges aufgestellt wird. Würfle nur einmal für alle
  Einheiten des Zuges in Reserve. Für Einheiten des Zuges, die in ein
  nicht-angeschlossenes Transportfahrzeug eingestiegen sind, werden die
  Reservewürfe separat durchgeführt.

\end{itemize}

\subsection{Inquisition}

\begin{itemize}

 \item Schnellschüsse durch den Träger des Psyocculum (oder seiner Einheit) auf
  eine Einheit die mindestens ein Modell mit der Sonderregel Psioniker,
  Psionischer Pilot oder einer Bruderschaft aus Psioniker/Hexern werden mit BF 1
  durchgeführt.

\end{itemize}

\subsection{Necrons}

\begin{itemize}

 \item Es werden keine Marker für Ewiges Leben oder Reanimationprotokoll
  platziert, sollte ein Modell überrannt werden. Marker für Ewiges Leben, die
  platziert wurden, bevor überrannt wurde, bleiben im Spiel.

 \item Ein erfolgreicher ``Verletzung ignorieren''-Wurf verhindert nicht den
  Rüstung ignorierenden Effekt durch Entropiestoß.

 \item Ein Kryptomen, Lord oder irgendein unabhängiges Charaktermodell, dass
  einer Einheit Eliminatoren beigetreten ist, profitiert von der Sonderregel
  ``Aus dem Nichts''.

 \item Der Todesstrahl der Todessichel ist eine rumpfmontierte Waffe, aber die
  zwei Punkte zum bestimmen der Trefferanzahl müssen nicht in Sicht zur Waffe
  sein.  Trotzdem können nur Modelle Wunden zugeteilt werden, die in Sicht und
  Reichweite sind. Einschließlich der Sichtlinie und Reichweite des
  Tesla-Annihilators.  Die Reichweite des Todesstrahls ist 12“ plus die 3W6“,
  die zur Bestimmung des zweiten Punktes geworfen werden.

 \item Der synchronisierte Tesla-Annihilator des Annihilatorgleiters ist eine
  rumpfmontierte Waffe.

 \item Die Geisterbarke kann die Fähigkeit ``Repair Barge'' nutzen um
  Kryptomanten/Lords, die einer Einheit Necronkrieger angeschlossen waren,
  wieder zu beleben. UCMs profitieren nicht von dieser Regel.

 \item Das Chronometrum darf nur für einzelne W6-Würfe genutzt werden (z.B.
  Trefferwürfe, Rüstungswürfe); nicht für Würfe mit mehreren W6 (z.B.
  Moralwerttest, Angriffsreichweite)

 \item Wenn ein Deathwing-Ritter durch Gedankenkontrollskarabäen übernommen
  wird, kann entweder der Necronspieler oder der DA-Spieler entscheiden, ob der
  Smite Mode aktiviert wird. Wenn ein Spieler bei einem höheren
  Initiative-Schritt zuschlägt, entscheidet dieser für die gesamte Einheit,
  welcher Modus verwendet wird..

 \item Vargard Obyrons Fähigkeit ``Mörderischer Gegenschlag'' findet nur
  Anwendung, wenn er entweder in einer Herausforderung kämpft oder nur ihm
  Wunden zugeordnet werden können.

\end{itemize}

\subsection{Orks}

\begin{itemize}

 \item Modelle, die durch Zogwarts Fluch zu einem Squig umgewandelt wurden,
  zählen nicht als Verlust. Erst wenn der Squig als Verlust entfernt wird, zählt
  das Originalmodell als Verlust.  Alle Sonderregeln, die ausgelöst werden, wenn
  das Originalmodell stirbt, werden ignoriert.

\end{itemize}

\subsection{Space Marines}

\begin{itemize}

 \item Ein Kontingent der Black Templars benutzt den entsprechenden Eintrag aus
  der Alliiertenmatrix des RB.

 \item Der W6 von Grav-Waffen gegen Fahrzeuge zählt in allen Belangen als
  Panzerungsdurchschlagswurf. Beispielsweise können Deckungswürfe durchgeführt
  werden.

 \item Grav-Waffen nutzen zum Verwunden den RW der Mehrheit der Einheit. Bei
  gleichstand wird der höhere (schlechtere) RW genommen.

 \item Grav-Waffen nutzen zum Verwunden den in der Phase aktuellen RW der
  Zieleinheit.

 \item Jede Ordensreliquie darf nur einmal pro Armee gewählt werden; unabhängig
  davon, wie viele Kontingente eingesetzt werden.

 \item Ein Modell darf nur mit einem Gegenstand von der Liste für
  Ordensreliquien ausgerüstet werden.

\end{itemize}

\subsection{Space Wolves}

\begin{itemize}

 \item Wenn ein Psioniker eines Zweckbündnisses in 24“-Umkreis zu einem
  Runenpriester eine Psi-Kraft wirkt, so muss der Space Wolves Spieler einen W6
  werfen um die Psi-Kraft eventuell zu bannen. Dieser Wurf wird anstelle des
  Psiresistenzwurfs durchgeführt.

 \item Sprunginfantrie und Sprung-Monströse-Kreaturen werden ganz normal vom
  ``Schlund des Weltenwolfes'' getroffen, egal ob sie im Flugmodus sind oder
  nicht.

\end{itemize}

\subsection{Sternenreich der Tau}

\begin{itemize}

 \item Ein Schildgenerator kostet 25 Punkte.

 \item Der Wurf für Aun'Vas Paradoxon of Dualität kann zusätzlich zu jedem
  anderen Schutzwurf durchgeführt werden. Waffen mit der allgemeinen Sonderregel
  ``Deckung ignorieren'' negieren nicht den Effekt des Paradoxon der Dualität..

 \item Eine Einheit, die von der Inkarnation der Elemente eines Himmlischen
  profitieren möchte, muss sich in 12“ zu ihm befinden, wenn sie dies nutzen
  möchte.  Beispielsweise muss im Falle des Anmut des Zephyrs die Rennbewegung
  innerhalb von 12“ zum Himmlischen enden um die Schnellschüsse abgeben zu
  können.

 \item Ein Geistteam muss sechs Modelle (Drohnen ausgenommen) umfassen, um den
  zweiten Fusionsblaster kaufen zu können.

 \item Ein Sonnenhai-Bomber beginnt das Spiel mit einer Pulsbombengenerator.

 \item Wenn eine Abfangdrohne aus einem Sonnenhai-Bomber aussteigt, hat die
  Bewegungsdistanz des Bombers keine Auswirkung auf die Schussattacke der
  Drohne.

 \item Schüsse eines Modell mit der Sonderregel Flugabwehr, dass eine Lenkrakete
  mit Hilfe eines Zielmarkierer und der Sonderregel ``Ziel erfasst'' auf einen
  Flieger oder eine fliegende monströse Kreatur im Flugmodus abfeuert, werden
  mit BF5 abgehandelt.
Schüsse von Modellen, die kein Flugabwehr haben, werden mit BF1 abgehandelt.

 \item Wenn ein Modell mit der Sonderregel Flugabwehr auf Flieger oder monströse
  fliegende Kreaturen im Flugmodus schießt und Lenkraketen mithilfe der
  Sonderregel ``Zielführung'' abfeuert, werden diese Schüsse BF5 abgehandelt.  Hat
  das Modell nicht Sonderregel Flugabwehr nicht, werden die Schüsse mit BF1
  abgehandelt.

 \item Ein Fahrzeug, dass nur Schnellschüsse abgeben kann, kann auch Lenkraketen
  mithilfe der Sonderregel ``Ziel erfasst'' nur mit BF1 abfeuern.

 \item Wenn mehrere Einheiten Abwehrfeuer geben und Zielmarkierer platziert
  werden, können diese durch eine andere Einheit, die später schießt genutzt
  werden.

 \item Für ein Punktverteidigungs-Zielrelais zählen Zielmarkierer als Waffen mit
  S5 oder weniger.

 \item Sowohl der Drohnensteuerung als auch das Gegenfeuer-Abwehrsystem
  Unterstützungssystem haben keine Auswirkung auf Drohnen, die Schnellschüsse
  abgeben.

 \item Multiple Zielerfassung kann im Abwehrfeuer nicht dazu benutzt werden, um
  eine Einheit zu beschießen, die nicht angreift.

\end{itemize}

\subsection{Tyraniden}

\begin{itemize}

 \item Ein Morgon, der die Fähigkeit Schrecken aus der Tiefe benutzt,
\begin{itemize}
 \item Der Morgon kann auch so schocken, dass er Modelle trifft, die im Nahkampf
  gebunden sind
 \item Wenn der Morgon einen Transporter zerstört, steigen die Insassen normal
  aus, bevor die 2. Schablone bzw der Morgon platziert wird. Dies kann dazu
  führen, dass der Morgon ein Schocktruppenmissgeschick erleidet.
 \item Wenn der Morgon platziert wird, ist die Mitte seines Bases über der Mitte
  der Schablone. Die Blickrichtung ist frei wählbar. Ein kleiner Teil des Bases
  wird über die Schablone hinausragen.
\end{itemize}

 \item Die Flügel der Schwarmdrude werden bei Überflugattacken ignoriert. Nutze
  den Kopf, Körper und das Base des Models.

 \item Die Schablone Geiferkanone der Schwarmdrude wird wie bei Flammenwaffen
  üblich am Base angelegt.

 \item Sporenminen, die durch Sporensalve erschaffen wurden, können in dieser
  Runde angreifen. Beachte, dass Sporenminen nicht angreifen dürfen, wenn sie
  geschockt sind.

 \item Die Entflammbar Sonderregel der Pyrovore betrifft nur Modelle in
  W6“-Umkreis.

\end{itemize}

